\begin{abstract}
Large Language Models (LLMs) are increasingly used by developers to generate Solidity smart contracts that manage billions of dollars in decentralized finance (DeFi). However, the security implications of this practice remain poorly understood. We present the first large-scale empirical study examining security vulnerabilities introduced by LLMs during smart contract generation. We generated 2,400 smart contracts across eight state-of-the-art LLMs---GPT-4o, GPT-5, Claude 3.5 Sonnet, Gemini 2.5 Flash-Lite, DeepSeek Coder V2, Qwen 2.5 Coder 32B, CodeLlama-34B, and Codestral 22B---spanning 20 contract categories including high-impact cross-chain bridge and flash loan contracts. Using a multi-tool analysis pipeline combining Slither, Mythril, and Semgrep, we detected and categorized 2,493 vulnerabilities in compilable contracts, mapped to 9 distinct CWE identifiers. Our results reveal that 54.7\% of compilable LLM-generated contracts contain at least one security vulnerability, with 26.8\% containing high-severity flaws that could lead to direct financial loss. Contrary to expectations, adversarial prompts exploiting gas optimization and simplicity biases do \emph{not} increase vulnerability rates, and this difference is not statistically significant ($p = 0.965$), suggesting that vulnerabilities are intrinsic to LLM code generation rather than prompt-dependent. We identify three LLM-specific vulnerability patterns---phantom guards, inherited vulnerabilities, and context window degradation---and find statistically significant differences across LLMs ($H = 132.65$, $p = 1.75 \times 10^{-25}$, Kruskal-Wallis test on compiled contracts). Compilation rates vary dramatically, from 91.0\% (GPT-5) to 34.7\% (CodeLlama), representing a critical and previously unreported quality dimension. When controlling for contract size, Codestral exhibits the highest vulnerability density (4.88 vulns/100 LOC), while GPT-5 achieves the lowest (0.99 vulns/100 LOC). Comparison with 241 verified Ethereum mainnet contracts reveals that vulnerability density is comparable (2.69 vulns/100 LOC for human vs.\ 0.99--4.88 for LLM), but human code benefits from professional auditing, monitoring, and defense-in-depth that LLM-generated code bypasses entirely. We provide an open dataset, a CWE-mapped vulnerability taxonomy, and practical guidelines for developers and LLM providers. Our findings underscore that the primary risk of LLM-generated smart contracts is not inferior code quality but the security infrastructure developers skip when trusting LLM output.
\end{abstract}

\begin{IEEEkeywords}
smart contracts, large language models, security vulnerabilities, Solidity, empirical study, adversarial prompts, static analysis, cross-chain bridges
\end{IEEEkeywords}
